\begin{center}\begin{tikzpicture}[font=\small,>=stealth,line join=round]
\draw[gray!40,rounded corners=3pt] (0,0) rectangle (14.2,8.3);\node at (7.1,7.700000000000001) {例4：基本段是指数曲线，翻折后仍隔4重复};
\begin{scope}[shift={(1.7923076923076926,3.8842857142857143)},x=1.153846153846154cm,y=1.3142857142857143cm]\draw[->,gray] (-0.6,0)--(4.6,0) node[right]{$x$};\draw[->,gray] (0,-1.7)--(0,1.8) node[above]{$y$};\draw[blue,thick,solid,domain=0:1,samples=180] plot(\x,{2^(\x)-1});\draw[blue,thick,solid,domain=1:2,samples=180] plot(\x,{2^(2-\x)-1});\draw[blue,thick,solid,domain=2:3,samples=180] plot(\x,{1-2^(\x-2)});\draw[blue,thick,solid,domain=3:4,samples=180] plot(\x,{1-2^(4-\x)});\fill[orange!85!black] (0,0) circle(2.2pt);\fill[orange!85!black] (1,1) circle(2.2pt);\fill[orange!85!black] (2,0) circle(2.2pt);\fill[orange!85!black] (3,-1) circle(2.2pt);\fill[orange!85!black] (4,0) circle(2.2pt);\foreach \v in {1,2,3,4}{\node[below=4pt] at (\v,-1.35){$\v$};}\end{scope}\node[anchor=west,align=left,text width=5.0cm] at (8.4,4.0) {$[0,1]$ 上 $f=2^x-1$\\[9pt]奇函数，轴 $x=1$\\[9pt]$f(1)=1,\ f(3)=-1$\\[9pt]$2023=505\times4+3$\\[9pt]$f(2023)=-1$};
\node[text width=12.6cm,align=center] at (7.1,.65) {仅画0到4的部分；将它按周期4延拓到全体实数，得到符合题设的函数。};
\end{tikzpicture}\end{center}
